\section{Relation thermique intrinsèque de l'échelle \(108\)}

La structure chronologique \(108\) fixe

\begin{equation}
k_B\Theta_{\rm clk}\log3
=\frac{h_{\rm int}}{108t_0}
=\frac{2\pi\hbar_{\rm int}}{108t_0}.
\label{eq:boltzmann-clock}
\end{equation}

Cette égalité ne détermine pas séparément une coordonnée numérique de \(k_B\) et une autre de \(\Theta_{\rm clk}\) sans normalisation supplémentaire. Elle détermine bien leur produit comme section d'énergie. Le facteur \(\log3\) a trois antécédents compatibles :

\begin{equation}
\log3=C_1+C_2-2C_0
=\frac{S_{\TRIT}}{k_B}
=\frac{E_P}{k_B\Theta_{\rm clk}}.
\label{eq:log3-triple}
\end{equation}

Les trois égalités correspondent aux réalisations résiduelle, informationnelle et
thermodynamique d'un même scalaire. La coïncidence de valeur n'identifie pas
leurs domaines respectifs.

\section{Formulations spectrales de la loi de déplacement de Wien}

Pour la densité spectrale par longueur d'onde, le maximum satisfait

\[
5(1-\ee^{-x_5})=x_5,
\qquad
x_5=5+W_0(-5\ee^{-5}).
\]

Avec \cref{eq:boltzmann-clock},

\begin{equation}
\lambda_{\max}(\Theta_{\rm clk})
=\frac{108\log3}{x_5}\,\ell_0.
\label{eq:wien-lambda}
\end{equation}

Pour la densité par fréquence,

\[
3(1-\ee^{-x_3})=x_3,
\qquad
x_3=3+W_0(-3\ee^{-3}),
\]

et

\begin{equation}
\nu_{\max}(\Theta_{\rm clk})
=\frac{x_3}{108t_0\log3}.
\label{eq:wien-frequency}
\end{equation}

Les maxima ne sont pas réciproques, puisque les densités spectrales se
transforment avec le jacobien correspondant. Le quotient \(x_5/x_3\)
appartient à la structure de la paramétrisation et ne représente pas une
incohérence.

\section{Loi de Stefan--Boltzmann et thermodynamique du gaz de photons}

Évaluer la loi de rayonnement en \(\Theta_{\rm clk}\) produit le flux intrinsèque

\begin{equation}
j_\star(\Theta_{\rm clk})
=\frac{2\pi^5}{15\,108^4(\log3)^4}\,
\frac{h_{\rm int}}{\ell_0^2t_0^2},
\label{eq:stefan-hmt}
\end{equation}

dans la normalisation déclarée. La constante de Stefan--Boltzmann n'est pas insérée comme primitive : elle résulte de l'intégration du spectre bosonique une fois l'action, la propagation et la température fixées.
Les lois spectrales employées sont celles de Planck et de Wien \cite{Planck,Wien} ; HMT fixe ici la section sur laquelle elles sont évaluées, et non leurs constantes par comparaison rétrospective.

En unités de l'échelle de Planck,

\begin{align}
n_\gamma\ell_P^3
&=\frac{2\zeta(3)}{\pi^2\log^3 3},
\label{eq:photon-number}\\
\frac{u_\gamma\ell_P^3}{E_P}
&=\frac{\pi^2}{15\log^4 3},
\label{eq:photon-energy}\\
\frac{s_\gamma\ell_P^3}{k_B}
&=\frac{4\pi^2}{45\log^3 3}.
\label{eq:photon-entropy}
\end{align}

\(\zeta(3)\) représente le moment bosonique d'ordre trois, tandis que
\(\log3\) détermine l'énergie de décision TRIT. L'équation conserve la
distinction entre les deux invariants.

La conductance thermique quantique s'exprime comme

\begin{equation}
\frac{\kappa_Qt_0}{k_B}=\frac{\pi^2}{324\log3}.
\label{eq:thermal-conductance}
\end{equation}

\section{Confluence métrologique des invariants thermiques et radiatifs}

\begin{longtable}{@{}p{0.18\textwidth}p{0.27\textwidth}p{0.45\textwidth}@{}}
\toprule
Objet & Équation interne & Significado estructural\\
\midrule
\endhead
\(R_K\) & \(Z_{\rm vac}/(2\alpha)\) & résistance quantique invariante bidirectionnelle\\
\(G_0\) & \(4\alpha/Z_{\rm vac}\) & conductancia complementaria\\
\(\Phi_0\) & \(\sqrt{hZ/(8\alpha)}\) & flux qui conserve l'orientation de l'action\\
\(K_J\) & \(\Phi_0^{-1}\) & relation réciproque entre fréquence et flux\\
\(E_P\) & \(k_B\Theta\log3\) & énergie du cycle et contenu informationnel TRIT\\
Wien & \(108\log3/x_5\) & extremum spectral à l'échelle chronologique\\
Apéry photonique & \(2\zeta(3)/(\pi^2\log^3 3)\) & densité de nombre bosonique\\
\bottomrule
\end{longtable}

\begin{remark}[Contrôle négatif]
Le bloc est réfuté si un changement de variable spectrale conserve
indûment le même extremum de Wien, si \(\zeta(3)\) remplace
\(\log3\), ou si des valeurs séparées sont attribuées à \(k_B\) et \(\Theta\) sans
spécifier de carte. Les identités
\cref{eq:quantum-electrical-identities,eq:planck-trit} fournissent des contrôles
algébriques directs.
\end{remark}

\subsection*{Portée logique et domaine de validité}
Quanta électriques et confluence Planck--TRIT : \emph{Théorème interne exact}. Wien, Stefan--Boltzmann et gaz photonique : \emph{Composition mathématique exacte} avec le continu bosonique déclaré. La carte SI reste une \emph{Comparaison métrologique externe}.

